\begin{fixedfigure}
\centering
\begin{tikzpicture}[>=Stealth,
    vertex/.style={circle,fill=black,inner sep=1.8pt},
    every node/.style={font=\small}]
\node[vertex,label=above left:\(x_0\)] (x0) at (0,1.7) {};
\node[vertex,label=above right:\(x_1\)] (x1) at (4.2,1.7) {};
\node[vertex,label=below left:\(y_0\)] (y0) at (0,0) {};
\node[vertex,label=below right:\(y_1\)] (y1) at (4.2,0) {};
\draw[->,thick] (x0) -- node[above] {\(h\)} (x1);
\draw[->,thick] (x0) -- node[left] {\(a_0\)} (y0);
\draw[->,thick] (x1) -- node[right] {\(a_1\)} (y1);
\draw[->,thick,dashed] (y0) -- node[above] {\(k\)} (y1);
\end{tikzpicture}
\caption{The solid arrows form the tree representing the joint-endpoint
lifting data \(B\).  Vertices denote \(n\)-cells and arrows denote
\((n+1)\)-cells; all vertices have the same lower boundary when \(n>0\).
Adjoining the dashed comparison \(k\) is one upper boundary attachment
\(I_{n+1}\), giving \(A_{n+1}^{01}\).  The square depicts boundary
data; its higher filling cells belong to the remaining upper completion,
which gives \(P_{n+1}\).}
\label{fig:joint-endpoints}
\end{fixedfigure}
